\chapter{Conclusions and future work}
\label{chapter:concl}

We framed the checking of an application repository against a regulation as grounded, selective classification:
every requirement receives a status, a confidence and citations, and the system may abstain in favour of a human. The
Compliance-as-Code AI Auditor implements this with a hybrid of deterministic policy-as-code checks and an LLM judge
(Claude via the Claude Code CLI in headless mode) constrained by retrieval, a strict citation contract, a Verifier and a
human-in-the-loop gate. It turns three policies (selected EU AI Act provisions, GDPR, an internal responsible-AI policy)
into 45 testable requirements and produces auditable reports in JSON, Markdown, SARIF and a web dashboard, with CI
regression gating and policy-version impact analysis.

On a benchmark of five hand-labelled AI applications (165 applicable labels), the main findings are:
\begin{itemize}[leftmargin=*]
  \item the hybrid mode is the most accurate (94.5\% accuracy, 86.6\% macro-F1), ahead of deterministic (82.4\%) and
        LLM-only (90.9\%);
  \item it needs 39\% fewer LLM calls than LLM-only, costs 41\% less and runs 45\% faster;
  \item it misses only 2 of 65 gaps, and all citations are valid;
  \item LLM-only mode is the most conservative (no false compliance), at the price of more false alarms.
\end{itemize}
Analysing the baseline errors improved accuracy by 3 (deterministic), 8 (hybrid) and 19 (LLM-only) points. The improvements were
comment-aware rules, graded rule confidence with a routing threshold chosen out-of-sample, and a rubric-aligned judge.
We recommend hybrid with Sonnet and $\tau=0.675$ as the default, LLM-only for high-stakes audits, and deterministic
mode for free, offline pre-commit checks.

\paragraph{Most important points.} (1) Rules and LLMs fail in complementary ways, so routing between them is worth more
than either alone; (2) the costliest auditor error, false compliance, can be controlled structurally (citations,
conflicts, abstention) rather than only statistically; (3) an auditor must be auditable: the record of who or what
decided each verdict, on which evidence, is as important as the verdict.

\paragraph{Future work.}
\begin{itemize}[leftmargin=*]
  \item \emph{Routing}: always route value-bearing checks (retention periods, amount thresholds) to the judge, which
        would remove the two remaining false compliances; learn the routing policy from review outcomes (active
        learning from the human queue).
  \item \emph{Requirements}: evaluate PolicyAnalyst-derived requirements against the curated ones; cover the full AI Act
        and harmonised standards; multilingual policies and repositories.
  \item \emph{Evidence}: dense or hybrid retrieval, cross-file reasoning (call graphs), runtime evidence (logs, metrics,
        model evaluations) in addition to static artefacts.
  \item \emph{Safety}: secret/PII redaction before prompting, local open-weight models for confidential code, calibration
        of the judge's confidence on review outcomes.
  \item \emph{Evaluation}: more and larger repositories labelled by independent annotators, with inter-annotator
        agreement.
\end{itemize}
